% CP-VI-0043
% target_chapter: VI/30 — Dimension of Schemes
% source_unit_id: IMP-DL-A5CEE6E541-P003
% source: Downloads/theory-of-algebraic-geometry-8.tex L748-L753
% solution_provenance: SOURCE_EXPLICIT_SOLUTION

\subsection*{Problem 3: Counterexamples for General Schemes}

\subsection*{Statement}

For each property (a), (c), and (d) in Question 2, find a scheme \(X\) for which the property does not hold.

\subsection*{Solution}

The statements in Problem 2 depend strongly on \(X\) being a variety, in particular on \(X\) being irreducible and equidimensional. For arbitrary schemes, the statements can fail.

We will use the same scheme to disprove all three properties.

Let
\[
X=\operatorname{Spec} k\ \sqcup\ \mathbb A^1_k.
\]
Equivalently,
\[
X=\operatorname{Spec} k\ \sqcup\ \operatorname{Spec} k[t].
\]

This is a scheme with two connected components:
\[
\operatorname{Spec} k
\]
and
\[
\mathbb A^1_k=\operatorname{Spec} k[t].
\]

Its dimension is
\[
\dim X=1,
\]
because
\[
\dim \operatorname{Spec} k=0
\]
and
\[
\dim \mathbb A^1_k=1.
\]
The dimension of a disjoint union is the supremum of the dimensions of the components, so
\[
\dim X=\max\{0,1\}=1.
\]

\subsection{Failure of Property (a)}

Property (a) says that for any closed point \(p\in X\),
\[
\dim X=\dim \mathcal O_{X,p}.
\]

Let \(p\) be the unique point of the component \(\operatorname{Spec} k\). Then \(p\) is a closed point of \(X\).

The local ring at \(p\) is
\[
\mathcal O_{X,p}\cong k.
\]
Therefore
\[
\dim \mathcal O_{X,p}=0.
\]

But
\[
\dim X=1.
\]

Hence
\[
\dim \mathcal O_{X,p}=0\neq 1=\dim X.
\]

Therefore property (a) fails for this scheme:
\[
\boxed{\dim X\neq \dim \mathcal O_{X,p}.}
\]

\subsection{Failure of Property (c)}

Property (c) says that if \(Y\subseteq X\) is closed, then
\[
\dim Y+\operatorname{codim}(Y,X)=\dim X.
\]

Let
\[
Y=\operatorname{Spec} k\subseteq X
\]
be the zero-dimensional connected component of \(X\). Since connected components of a disjoint union are both open and closed, \(Y\) is closed in \(X\).

We have
\[
\dim Y=0.
\]

The only irreducible closed subset of \(Y\) is \(Y\) itself. Its generic point is the unique point \(p\in \operatorname{Spec} k\).

The codimension of \(Y\) in \(X\) is
\[
\operatorname{codim}(Y,X)=\dim \mathcal O_{X,p}.
\]
But as above,
\[
\mathcal O_{X,p}\cong k,
\]
so
\[
\dim \mathcal O_{X,p}=0.
\]
Therefore
\[
\operatorname{codim}(Y,X)=0.
\]

Thus
\[
\dim Y+\operatorname{codim}(Y,X)=0+0=0.
\]

But
\[
\dim X=1.
\]

Hence
\[
\dim Y+\operatorname{codim}(Y,X)=0\neq 1=\dim X.
\]

Therefore property (c) fails:
\[
\boxed{\dim Y+\operatorname{codim}(Y,X)\neq \dim X.}
\]

\subsection{Failure of Property (d)}

Property (d) says that if \(U\subseteq X\) is a nonempty open subset, then
\[
\dim U=\dim X.
\]

Again take
\[
U=\operatorname{Spec} k\subseteq X.
\]
Since \(X\) is a disjoint union,
\[
\operatorname{Spec} k
\]
is both open and closed in \(X\). Thus \(U\) is a nonempty open subset of \(X\).

But
\[
\dim U=0,
\]
whereas
\[
\dim X=1.
\]

Therefore
\[
\dim U=0\neq 1=\dim X.
\]

Hence property (d) fails:
\[
\boxed{\dim U\neq \dim X.}
\]

\subsection*{Conclusion}

The scheme
\[
X=\operatorname{Spec} k\sqcup \mathbb A^1_k
\]
shows that properties (a), (c), and (d) from Problem 2 do not hold for arbitrary schemes.

The reason is that this scheme is not irreducible and not equidimensional. It has one component of dimension \(0\) and another component of dimension \(1\). Therefore closed points and open subsets lying entirely in the zero-dimensional component do not detect the full dimension of \(X\).
